\section{An eleven-positive equation of degree 213}
\label{sec:strong-pell-budget}

This companion lowers the incoming centered eleven-positive equation from
degree $257$ and $114$ operations to degree $213$ and $112$ operations.
The incoming degree-$277$ equation still has one fewer operation
\cite{joint}.  The new mechanism is to use growth forced by the strong
Pell equation to bound the transport guard, then center the norm at
\emph{both} the strong and auxiliary roots.  All statements in this
section concern ordinary proofs formalizable in ZFC.  The compact44
native compiler and the normalized source calibration are inherited;
this is not a new end-to-end proof-assistant certificate.

\subsection{The complete polynomial}

The unknowns are
\[
 (j,F,\alpha,i,T,X,Y,n,\eta,\zeta,\gamma)\in\Zp^{11}.
\]
Use the entire compact44 coefficient tuple
$(\beta,K_0,M_C,M_F)$, with $\beta=2^{590069}-1$ and the already shifted
$M_F$.  The companion includes their actual integer payloads and binding;
these ports are not interchangeable with a different native tuple.
Every following name abbreviates a polynomial and adds no unknown:
\begin{align*}
 J&=3j+1,&q&=\beta J+1,&C&=q-F-\alpha,&E&=XY,\\
 k&=\eta+\zeta,&c&=kY+\eta,&a&=E+Y,&H&=4a+3,\\
 \Delta&=a^2+H,&\ell&=EkY,&I&=ic^2,&D_f&=1+\Delta I^2,\\
 Q&=\Delta^2I^2,&\delta&=Q-1,&M&=q^3(q-1).
\end{align*}
Further put
\begin{align*}
 R&=((q-1)(q-F)+\alpha)(q^2-1)+(M_C+qM_F)J,\\
 B_t&=(qK_0+X)C-qF,&D_\tau&=\ell(\ell+k)+1,\\
 V&=D_f(cT-R)-c,&\bar f&=(a+2)I,\\
 h_*&=i^3+cT,&W&=1+(q-1)h_*,\\
 \bar Z&=M(\ell+k+n+cT+V+\bar f)+(q-1)Y-q^3WB_t.
\end{align*}
The unchanged main carrier is
\begin{align*}
 v&=k-R-1,&S&=4qY,&U&=qS,\\
 D_*&=U(X+ac)+H(U\gamma+SX-v),\\
 A_1&=D_*^2-U^2(1+\Delta c^2).
\end{align*}
The octic norm is evaluated by
\begin{align*}
 A_r&=M^2D_\tau,&B_r&=M^2D_f,&C_r&=M^2(QV^2-1),\\
 T_N&=\delta^2A_rB_r,&R_N&=\delta C_r,\\
 J_N&=\delta(\bar Z^2-A_r-B_r)-C_r,\\
 N_c&=J_N^2-4\{T_N+R_N(A_r+B_r)\},\\
 N_8&=N_c^2-64T_NR_N\bar Z^2.
\end{align*}
The single equation is
\begin{equation}
 \boxed{\quad P_{11}=N_8+A_1+N_8A_1=0.\quad}
 \label{eq:pell-budget11}
\end{equation}
It uses only fixed integer coefficients, fixed exponents, addition,
subtraction, and multiplication.  Since $4\mid U$, both $A_1$ and $N_8$
are squares modulo $4$.  Thus $(1+N_8)(1+A_1)=1$ forces both factors to
be $1$: the alternative $(-1,-1)$ would require residues $-2$.
Consequently \eqref{eq:pell-budget11} is exactly the conjunction
$N_8=A_1=0$ on its stated domain.

\subsection{Source boundary and recovery before positivity}

The inherited weighted source theorem \cite{source70,joint} concerns
positive coordinates satisfying the main residual, the strong norm,
the auxiliary norm, and the first carrier
\begin{align*}
 f^2-\Delta(ic^2)^2&=1,\\
 QV^2-(Q-1)(V+u)^2&=1,\\
 Z^2&=M^2D_\tau,\\
 Z&=M(\ell+k+r)+(q-1)Y-q^3WB_t.
\end{align*}
Its soundness holds for every positive integer $h_*$.  For canonical
witnesses, after the strong and auxiliary pairs have been chosen, a
chosen weight can be lifted by
\[
 r=Wt_0+\tau-\ell-k-s>qW-2k,
 \qquad t_0\ge q,\quad 0<s<k,\quad \ell<\tau<\ell+k.
\]
The source has a positive zero precisely when the fixed ZFC calculus has
a finite refutation.  The authentic mask bounds are
\[
 K_0\ge1,\quad 0<M_C<\beta,\quad0<M_F<2\beta,
 \qquad 0<(M_C+qM_F)J<2q^2.
\]
Once $C>0$, these imply $3q<R<q^4-2$.

Here is the part of the norm recovery that must precede $C>0$.
For every positive input, $I\ge c^2$, $\Delta\ge15$, and
\[
 1<D_\tau<D_f<Q,\qquad |V|>1.
\]
The rational radicand
$D_y=(QV^2-1)/(Q-1)>V^2$ is distinct from $D_\tau,D_f$.
Moreover, with $t=\sqrt{D_\tau}$, $f=\sqrt{D_f}$,
$y=\sqrt{D_y}$, no signed sum of $t,f,y$ is zero.
If $cT-R=0$, then $f>t+y$; otherwise $y>f+t$.
These inequalities depend only on $I\ge c^2$, not on $C$ or a bound on
$R$; their elementary proof is retained in the companion source audit.

The universal three-root norm therefore gives, at $N_8=0$,
\begin{equation}
 \bar Z/M=\tau_*+f_*+y_*,\qquad
 \tau_*^2=D_\tau,\quad f_*^2=D_f,\quad y_*^2=D_y.
 \label{eq:budget-radical-recovery}
\end{equation}
All three roots are rational.  One can prove this without assuming
independent square classes: for signed roots $x_1,x_2,x_3$ with rational
nonzero sum $t_0$, put
\[
 u_0=(t_0^2-x_1^2-x_2^2-x_3^2)/2,
 \qquad
 v_0=\frac{u_0^2-x_1^2x_2^2-x_1^2x_3^2-x_2^2x_3^2}{2t_0}.
\]
Then $u_0$ and $v_0$ are rational and
$x_1=x_1^2(x_1^2+u_0)/(x_1^2t_0+v_0)$.
The denominator is $x_1(x_1+x_2)(x_1+x_3)\ne0$, by distinctness of
the radicands; use the cyclic formulas as well.
Thus $\tau_*,f_*$ are integers.  The adjacent-coefficient Pell descent
from $QV^2-(Q-1)y_*^2=1$ makes $y_*$ an integer too
\cite{joint}.  In particular $M\mid\bar Z$.

Reduction modulo $q^3$ and $q-1$ now gives $q^3\mid Y$ and
$q-1\mid B_t$, independently of the sign of $C$.
The main residual makes $g=|D_*/U|$ an integer and
$g^2-\Delta c^2=1$.  Writing
\[
 \chi_A(r)+\psi_A(r)\sqrt{A^2-1}
       =(A+\sqrt{A^2-1})^r,
\]
Pell classification gives $c=\psi_{a+2}(p)$ and
$k=2\psi_{2XY^2+1}(m_1)$.  Here $p\ge2$ and $m_1\ge2$:
$m_1=1$ would give $k=2$, $c=2Y+1<2(a+2)$.
Consequently, still before proving $C>0$,
\begin{equation}
 c>kY\ge4(2XY^2+1)Y>8XY^3\ge8q^9.
 \label{eq:budget-c-growth}
\end{equation}

\subsection{The new conditional growth bound}

\begin{lemma}[Strong-Pell growth budget]
Suppose $q\ge2$, $Y\ge q^3$, $A\ge2Y+2$,
$\Delta=A^2-1$, $c=\psi_A(p)\ge6$, and positive integers $f,i$
satisfy $f^2-\Delta(ic^2)^2=1$.  Then
\[
 i>8q\Delta c^4,\qquad i^3>6qf^2,\qquad i^3>q^2,
 \qquad Aic^2<i^2.
\]
If $0<k<c$, also $i^3>Aic^2+k$.
\end{lemma}
\begin{proof}
Classify $f=\chi_A(m)$ and $ic^2=\psi_A(m)$.
Strong divisibility gives $p\mid m$: the addition law and
$\gcd(\chi_A(r),\psi_A(r))=1$ give the Euclidean identity
$\gcd(\psi_A(r),\psi_A(s))=\psi_A(\gcd(r,s))$.
Write $m=pv$ and $g=\chi_A(p)$.  Composition gives
$ic=\psi_g(v)$.  As $g^2\equiv1\pmod c$, recurrence gives
$\psi_g(v)\equiv vg^{v-1}\pmod c$; hence $c\mid v$.
Thus $v\ge c\ge6$ and $ic>g^5$.  Since
$g^2=1+\Delta c^2$ and $\Delta>4Y^2$,
\[
 i>\Delta^{5/2}c^4
   >8Y^3\Delta c^4\ge8q^9\Delta c^4\ge8q\Delta c^4.
\]
Putting $H_0=\Delta c^4\ge15$ gives
$i^3>8qH_0i^2>6q(H_0i^2+1)=6qf^2$.
Also $i>q$, $i>Ac^2$, and $i>c>k$.
Therefore $Aic^2<i^2$ and $i^3-i^2>k$.
\end{proof}

Apply the lemma after \eqref{eq:budget-radical-recovery}, with
$A=a+2$ and $f^2=D_f$.  It gives
\begin{equation}
 i^3>6qD_f,\qquad i^3>q^2,\qquad i^3>\bar f+k.
 \label{eq:budget-bounds}
\end{equation}
Also $\bar f>f$, since $\bar f^2-D_f=I^2-1>0$.
These are consequences of the recovered Pell equations; no assertion
of these bounds on arbitrary positive inputs is needed.

\subsection{The noncircular positivity bootstrap}

Suppose $C\le0$, put $d=-C$ and $B=-B_t>0$.  Then
\[
 B=(qK_0+X)d+qF,\qquad
 R=[q(q-F)+d](q^2-1)+(M_C+qM_F)J.
\]
The mask bounds imply
\begin{equation}
 R<qB+q^4+q^2,\qquad R<0\Longrightarrow -R<q^2B.
 \label{eq:budget-rank-bounds}
\end{equation}
Let $s=Y/q^3>0$.  Equation \eqref{eq:budget-radical-recovery} becomes
\[
 \tau_*+f_*+y_*
  =\ell+k+n+cT+V+\bar f+s+\frac{WB}{q-1}.
\]
Its left side is at most $t+f+y$, with $t<\ell+k$.
We derive a contradiction in both possible signs of $V$.

As $\delta>D_f$,
\[
 0<y-|V|=\frac{V^2-1}{\delta(y+|V|)}<\frac{|V|}{\delta}.
\]
If $V>0$, the integer $b=cT-R$ is positive, so $y-V<b$ and
\[
 cT+V+\bar f-f-y>cT-b=R.
\]
For $R\ge0$ the right side of the radical equality already exceeds
$t+f+y$.  For $R<0$, the remaining term exceeds
$h_*B\ge i^3B>q^2B>-R$, with the same contradiction.

If $V<0$, then $R\ge cT\ge c>8q^9$.  Thus
$q^4+q^2<R/2$ and \eqref{eq:budget-rank-bounds} gives $R<2qB$.
We have $|V|<D_fR$ and $y-|V|<R$, whence
\[
 cT+V+\bar f-f-y>-(2D_f+1)R>-3D_fR.
\]
But $WB/(q-1)>h_*B\ge i^3B>6qD_fB>3D_fR$, again a
contradiction.  The case $V=0$ was excluded earlier.  Therefore $C>0$.

Now $3q<R<q^4-2$, and \eqref{eq:budget-c-growth} gives $c>2R$.
Hence $V>1$ and the positive-gap estimate strengthens to
\[
 cT+V+\bar f-f-y>R>0.
\]
Restore positive source witnesses by
\[
 f_{\rm old}=f,\qquad u=y-V,\qquad
 r=n+cT+V+\bar f-f_*-y_*.
\]
Here $u>0$ and $r>n>0$ for every permitted choice of signs.
Subtracting $M(f_*+y_*)$ from $\bar Z$ gives exactly the first source
carrier $Z=M\tau_*$.  All four source equations hold, so source
soundness yields a finite ZFC refutation.

Conversely, choose canonical outer, strong, and auxiliary source
witnesses first, then set $h_*=i^3+cT$.
By \eqref{eq:budget-bounds}, $h_*>cT+\bar f+k$; the weighted
source construction supplies $r>qW-2k>cT+\bar f$.
Take
\[
 n=r-cT-\bar f+f+u>0.
\]
Then $\bar Z=M(\tau+f+y)$, so $N_8=0$, and the unchanged canonical
main carrier gives $A_1=0$.  No new congruence on the strong pair is
required.  The witness choices have no dependence loop.
Composing both implications with the inherited source calibration proves
\[
 \mathrm{ZFC}\vdash
 \bigl(\nexists\boldsymbol{x}\in\Zp^{11}\;
       P_{11}(\boldsymbol{x})=0\bigr)\ \longleftrightarrow\ \ConZFC.
\]

\subsection{Exact degree and reproducible finite evidence}

Put $e=I^2-1$ and $z=\bar Z-M(V+\bar f)$, so
$D_f=\bar f^2-e$ and $\deg z=9$.
Expand the universal norm using
\[
 x=M(V+t)+z,\quad A=M^2D,\quad B=M^2(t^2-e),\quad
 C_{\rm num}=M^2(QV^2-1),\quad\delta=Q-1.
\]
With weights $(18,18,4,7,9,10,8)$ on $(Q,V,M,t,z,e,D)$,
the exact $1110$-term integer expansion has degree $200$ and leading form
\begin{equation}
 64M^6Q^4V^4t^2z^2.
 \label{eq:budget-leading-norm}
\end{equation}
The checker expands this identity exactly; naive gate-degree propagation
does not suffice.  Write $q_0=3\beta j$, $k_0=\eta+\zeta$, and
\[
 B_{t,0}=(K_0q_0+X)(q_0-F-\alpha)-q_0F.
\]
The relevant leading forms are
\begin{align*}
 D_{f,0}&=i^2k_0^4X^2Y^6,&Q_0&=i^2k_0^4X^4Y^8,\\
 R_0&=q_0^3(q_0-F),&V_0&=-D_{f,0}R_0,\\
 \bar f_0&=ik_0^2XY^3,&M_0&=q_0^4,\\
 z_0&=-q_0^4(i^3+k_0YT)B_{t,0},\\
 (A_1)_{13}&=32q_0^2k_0X^2Y^4
       \{4q_0^2Y\gamma+q_0^3(q_0-F)\}.
\end{align*}
Therefore $\deg P_{11}\le213$.  Its coefficient at
\[
 j^{48}i^{24}\eta^{37}X^{30}Y^{67}F^4\alpha^2\gamma
\]
is exactly $2^{13}(3\beta)^{48}>0$.  The minimum $q_0$ degree and
the maximum displayed $i$ degree select this coefficient uniquely in
\eqref{eq:budget-leading-norm} times $(A_1)_{13}$; it is independent of
$K_0,M_C,M_F$.  Thus the exact degree on the actual numeral tuple is
$213$.

The standard-library checker passes normally and with Python assertions
disabled.  It checks the complete prefix, the octic conjugate identity,
the weighted leading form, the nonzero degree certificate, coefficient
hashes, deterministic rebuild, all eleven live inputs, and all $112$
live arithmetic operations.  Constructed Pell examples supplement the
unbounded proof above; they do not test consistency or replace the
inherited semantic calibration.
